\documentclass[11pt]{article}
\usepackage{amsmath,amssymb}
\title{The offset in a pose reading also shifts the remaining curve}
\date{26 November 2026}
\begin{document}
\maketitle

\paragraph{Interface.}
In one corridor-face coordinate, let $q(t)$ be the issued quadratic
reference, $0\leq t\leq1$. At a reading cut the actual path parameter
$s$ is unknown in a warranted prior interval $[l,r]$. The available
observation says that $q(s)+z\in[L,U]$, where the calibration offset
$z$ is initially in $[a,b]$. Motion after the reading follows a
parameter in $[s,1]$. The verification envelope takes all these
parameters seriously. During the future cut the offset may be any
$z'\in[a,b]$ with $|z'-z|\leq d$; the static law has $d=0$.
For $d>0$ the envelope allows time to elapse while the path
parameter stays at $s$. This is a bounded-difference law with all
such pairs admitted, not a drift-rate theorem.
The ceiling $c$ is inclusive. Tracking error, other faces, approach
dynamics, logging delays, and how the motion monitor warrants $[l,r]$
are outside this scalar model. The distinction at issue is that $z$
is the \emph{same} offset used to obtain the reading and start the
future error bound. Giving each use an independent $[a,b]$ lowers
what can be certified.

\paragraph{Exact calculation.}
For any allowed subinterval $J$ of past parameters, the range of
$q$ is the closed interval $[m_J,M_J]$; continuity yields every
ordinate between its extremizers. Intersect this with the reading
constraint,
\[
  [A_J,B_J]=[m_J,M_J]\cap[L-b,U-a].
\]
An empty intersection for $J=[l,r]$ means inconsistent evidence,
not successful clearance. For nonempty intersections, the largest
feasible future offset given observed-path ordinate $y=q(s)$ is
$\min(b,U-y+d)$. Indeed take the greatest admissible current
offset $\min(b,U-y)$; the evidence constraint is feasible at that
offset by the interval test, and it can rise by $d$ or hit $b$.

At any $s$ a quadratic peaks on $[s,1]$ at $s$, at $1$, or at
its fixed vertex $v$ if concave and $v$ is within the suffix.
For the first case, optimize the nondecreasing function
$\min(y+b,U+d)$ over $[A_{[l,r]},B_{[l,r]}]$, attaining its right endpoint.
For a fixed future parameter $f\in\{1,v\}$, use
$J=[l,\min(r,f)]$, omitting empty subintervals and intersections.
The greatest feasible future height is
\[
    q(f)+\min(b,U-A_J+d).
\]
The minimum feasible $y$ maximizes this formula. In each case the
chosen ordinate is attained by some past parameter in $J$, and
the chosen offset pair and future parameter are jointly feasible.
Take the largest of the resulting numbers. This proves exact
classification and attainment rather than a separately upper-bounded
sum. The comparison uses a separate quadratic root-sign oracle to
decide the existence of violations, and 900 exact instances with
and without allowed drift. The prior checked interval lemma in the
parent provides the same extremizers, but the conditioning theorem
here is a separate mathematical proof, not a new Lean file.

\paragraph{Sharp distinction in a curved case.}
Take $q(s)=4s(1-s)$, $s\in[1/3,1]$, $z\in[0,1/3]$, and
$[L,U]=[9/10,1]$. For $1/3\leq s\leq1/2$, the future peak of
$q$ is $1$, but $q(s)\geq8/9$ forces $z\leq1/9$. The pair
$s=1/3,z=1/9$ reaches height $10/9$ at the peak if $d=0$.
For $s\geq1/2$, $q$ decreases, so future physical height is
at most $1$ with a static offset. The marginal bounds alone are
exactly $\max q=1$ and $\max z=1/3$: the latter is compatible with,
for example, the late descending branch and $q(s)=2/3$. The maxima
do not coexist. The joint bound admits ceiling $10/9$ when the
sum of separate maxima $4/3$ would refuse.

The exact bound allowing drift is
$1+\min(1/3,1/9+d)$: the early branch attains it and the late
branch is bounded by $\min(4/3,1+d)$.
At this ceiling any positive drift can invalidate the decision;
after $d=2/9$ the joint advantage over the marginal sum disappears.
If the asserted prior lower bound is not warranted, an earlier
$s$ on the ascending branch can agree with the same reading and
violate the ceiling even without drift. The adjacent source replay
does not warrant this reading or progress interface. An equally
informed planner retaining the offset relation can use the same
calculation. This is an exact scalar conditioning result, not a
physical clearance, learned drift or contract acceptance claim.

\end{document}
